% ======================================================================
%  MACROS — Sistema de design editorial do livro (A4, marinho + ouro + papel)
% ======================================================================
\usepackage[T1]{fontenc}
\usepackage[utf8]{inputenc}
\usepackage[brazil]{babel}
\usepackage{lmodern}
\renewcommand{\familydefault}{\sfdefault}
\usepackage{microtype}
\usepackage{textcomp}
\usepackage[table]{xcolor}
\usepackage[a4paper,inner=2.65cm,outer=2.25cm,bottom=1.8cm,top=2.7cm,headheight=22pt,headsep=0.2cm,footskip=12pt]{geometry}
\usepackage{booktabs,colortbl,tabularx,array,ragged2e,longtable,multirow}
\usepackage{enumitem}
\usepackage{chngcntr}
\usepackage{titlesec}
\usepackage{tocloft}
\usepackage{fancyhdr}
\usepackage{graphicx}
\usepackage{fancyvrb}
\usepackage{tcolorbox}
\tcbuselibrary{breakable,skins,listings}
\usepackage{tikz}
\usetikzlibrary{positioning,arrows.meta,shapes.geometric,calc,matrix,fit,backgrounds}
\usepackage{hyperref}

% ----------------------------------------------------------------------
%  PALETA
% ----------------------------------------------------------------------
\definecolor{navy}{RGB}{48,17,76}          % roxo profundo do cabeçalho/rodapé de referência
\definecolor{steel}{RGB}{187,125,242}      % violeta luminoso de destaque
\definecolor{steelV}{RGB}{48,17,76}        % roxo profundo para fundos
\definecolor{teal}{RGB}{187,125,242}       % acento unificado
\definecolor{violet}{RGB}{187,125,242}     % destaque violeta
\definecolor{amber}{RGB}{187,125,242}      % acento unificado
\definecolor{green}{RGB}{187,125,242}      % acento unificado
\definecolor{red}{RGB}{187,125,242}        % acento unificado
\definecolor{slate}{RGB}{170,170,170}      % cinza secundário
\definecolor{papel}{RGB}{34,34,34}         % fundo carvão
\definecolor{gold}{RGB}{187,125,242}       % compatibilidade: detalhes em violeta
\definecolor{mist}{RGB}{52,52,52}          % painel cinza
\definecolor{ink}{RGB}{34,34,34}           % fundo carvão da referência
\definecolor{panel}{RGB}{42,42,42}         % painel principal
\definecolor{panelAlt}{RGB}{48,48,48}      % painel alternado
\definecolor{textmain}{RGB}{204,204,204}   % texto cinza-claro
\definecolor{textmuted}{RGB}{178,178,178}  % texto secundário
\definecolor{lineDark}{RGB}{92,92,92}      % divisórias
\colorlet{plumDeep}{navy}
\colorlet{plumAccent}{steel}
\colorlet{.}{textmain} % padrão do documento: evita resets para preto no tema escuro

% ----------------------------------------------------------------------
%  COLUNAS / TABELAS
% ----------------------------------------------------------------------
\newcolumntype{Y}{>{\RaggedRight\arraybackslash}X}
\newcolumntype{L}{>{\RaggedRight\arraybackslash}p{3.2cm}}
\setlength{\tabcolsep}{4pt}
\arrayrulecolor{lineDark}

% ----------------------------------------------------------------------
%  TÍTULOS (part / chapter / section)
% ----------------------------------------------------------------------
\titleformat{\part}[display]
  {\thispagestyle{plain}\sffamily\bfseries\raggedright}
  {\fontsize{90}{90}\selectfont{\color{gold!35}\thepart}\\[-24pt]{\color{gold}\rule{\linewidth}{2.5pt}}}
  {10pt}{\fontsize{26}{30}\selectfont\color{textmain}}
  [{\vspace{6pt}{\color{steel}\hrule width \linewidth height 1pt}\vspace{4pt}\color{textmuted}\normalsize\normalfont\itshape\partcontent}]

\titleformat{\chapter}[display]
  {\thispagestyle{fancy}\sffamily\bfseries\filleft}
  {\noindent\color{plumAccent}\small\MakeUppercase{\chapterlabel}\hfill
   {\fontsize{48}{48}\selectfont\thechapter\hspace{2mm}}\colorbox{plumAccent}{\rule{0.08\linewidth}{0.9cm}}}
  {4pt}{\fontsize{25}{29}\selectfont\color{textmain}}
  [\vspace{5pt}{\color{plumAccent}\rule{\linewidth}{1.8pt}}]

\titleformat{\section}
  {\normalfont\sffamily\bfseries\large\color{textmain}}
  {\thesection.}{0.6em}{}
  [\vspace{1pt}{\color{lineDark}\rule{0.35\linewidth}{0.6pt}}]

\titleformat{\subsection}
  {\normalfont\sffamily\bfseries\small\color{gold}}
  {\thesubsection.}{0.6em}{}

% Babel pt-BR substitui "Capítulo" por "Capítulo" (default) — garantimos os rótulos:
\renewcommand\chaptername{Elemento}
\newcommand\chapterlabel{Elemento}
\newcommand\partcontent{Parte do manual técnico-dev-científico do OpenCode Ecosystem Core}

% ----------------------------------------------------------------------
%  CABEÇALHO E RODAPÉ
% ----------------------------------------------------------------------
% Marcas correntes: título da seção no alto e capítulo no rodapé.
\renewcommand{\chaptermark}[1]{\markboth{Capítulo \thechapter. #1}{}}
\renewcommand{\sectionmark}[1]{\markright{#1}}
\pagestyle{fancy}
\fancyhf{}
% Texto alinhado junto às bordas da página e centralizado verticalmente nas faixas.
\fancyhfoffset[L]{2.25cm}
\fancyhfoffset[R]{1.85cm}
\fancyhead[L]{\raisebox{31pt}[0pt][0pt]{\parbox[b]{0.92\headwidth}{\raggedright\sffamily\bfseries\footnotesize\color{white}\MakeUppercase{\rightmark}}}}
\fancyhead[R]{}
\fancyfoot[L]{\raisebox{-20pt}[0pt][0pt]{\parbox[b]{0.88\headwidth}{\raggedright\sffamily\bfseries\footnotesize\color{white}\MakeUppercase{\leftmark}}}}
\fancyfoot[R]{\raisebox{-20pt}[0pt][0pt]{\footnotesize\sffamily\bfseries\color{white}\thepage}}
\renewcommand{\headrulewidth}{0pt}
\setlength{\headheight}{22pt}
% Faixas roxas contínuas, inspiradas no cabeçalho e rodapé do livro de referência.
\AddToHook{shipout/background}{%
  \begin{tikzpicture}[remember picture,overlay]
    \fill[plumDeep] (current page.north west) rectangle ([yshift=-24mm]current page.north east);
    \fill[plumDeep] (current page.south west) rectangle ([yshift=15mm]current page.south east);
  \end{tikzpicture}%
}
\fancypagestyle{plain}{%
  \fancyhf{}
  \fancyhead[L]{\raisebox{31pt}[0pt][0pt]{\parbox[b]{0.92\headwidth}{\raggedright\sffamily\bfseries\footnotesize\color{white}\MakeUppercase{\rightmark}}}}
  \fancyfoot[L]{\raisebox{-20pt}[0pt][0pt]{\parbox[b]{0.88\headwidth}{\raggedright\sffamily\bfseries\footnotesize\color{white}\MakeUppercase{\leftmark}}}}
  \fancyfoot[R]{\raisebox{-20pt}[0pt][0pt]{\footnotesize\sffamily\bfseries\color{white}\thepage}}
  \renewcommand{\headrulewidth}{0pt}
}

% ----------------------------------------------------------------------
%  Cabeçalho de ELEMENTO — painel preto com título branco
% ----------------------------------------------------------------------
\newcommand{\elemento}[2]{%
  \par\medskip
  \begin{tcolorbox}[enhanced,breakable,boxrule=0.8pt,arc=1.6mm,
    colback=plumDeep,colbacktitle=plumDeep,colframe=plumDeep,coltitle=white,
    borderline west={1.5pt}{0pt}{gold},
    title={\hspace{2mm}\parbox[b]{\dimexpr\linewidth-4mm}{#1}},
    fonttitle=\bfseries\normalsize,
    left=2mm,right=2mm,top=1.6mm,bottom=2.2mm]
    \footnotesize\color{white}#2
  \end{tcolorbox}
  \par\nointerlineskip\vspace{1mm}\par}

\newcommand{\metainfo}[2]{\footnotesize\color{white}\textbf{#1:}~#2}

% ----------------------------------------------------------------------
%  CITAÇÃO AUDITADA — nota de rodapé verificável (regra R110 / anti-overclaim)
%  args: {referência}{DOI ativo}{relevância ao contexto (pt-BR)}
%        {trecho original}{tradução pt-BR do trecho}
%  Uso SOMENTE em parágrafo corrido (fora de tcolorbox) e com DOI
%  previamente verificado por sondagem/webfetch — nunca de memória.
% ----------------------------------------------------------------------
% ----------------------------------------------------------------------
%  CITAÇÃO AUDITADA — nota de rodapé verificável (regra R110 / anti-overclaim)
%  Sistema NUMÉRICO de chamada em nota de rodapé, conforme ABNT NBR
%  10520:2023; a referência completa segue a ABNT NBR 6023:2018.
%  args: {referência ABNT (sem DOI)}{DOI ativo}{relevância ao contexto (pt-BR)}
%        {trecho original}{tradução pt-BR do trecho}{AUTOR, ANO, p. PÁGINA}
%  Uso SOMENTE em parágrafo corrido (fora de tcolorbox) e com DOI
%  previamente verificado por sondagem/webfetch — nunca de memória.
% ----------------------------------------------------------------------
\makeatletter
\newcommand{\citref}[6]{%
  % Usa o contador global, inclusive dentro de caixas tcolorbox.
  \begingroup
  \def\@mpfn{footnote}\def\thempfn{\thefootnote}%
  (#6)\footnote{\footnotesize\color{white}\textbf{Referência bibliográfica:} #1.%
    \par\smallskip\textbf{DOI:} \href{https://doi.org/#2}{\texttt{#2}}.%
    \par\smallskip\textbf{Relevância a este contexto:} #3.%
    \par\smallskip\textbf{Citação direta (ABNT NBR 10520:2023):} ``{\itshape #4}''.%
    \par\smallskip\textbf{Tradução livre (pt-BR):} ``#5''.}%
  \endgroup
}
\makeatother

% Notas sempre claras no tema escuro, inclusive quando vêm de caixas ou referências.
% A numeração bibliográfica segue sem reiniciar nos capítulos.
\counterwithout{footnote}{chapter}
\renewcommand{\thefootnote}{\arabic{footnote}}
\makeatletter
\long\def\@makefntext#1{%
  \parindent 1em\noindent
  {\everypar{\color{white}}\color{white}%
    \hb@xt@1.8em{\hss\@makefnmark}#1}%
}
\renewcommand{\footnoterule}{%
  \kern-3pt
  {\color{white}\hrule width .4\columnwidth height .4pt}%
  \kern2.6pt
}
% Evita que a continuação de uma nota longa perca a cor ao atravessar páginas.
\interfootnotelinepenalty=10000
\makeatother

% ----------------------------------------------------------------------
%  BIBLIOTECA DE CAMPOS DA FICHA (boxes temáticos)
% ----------------------------------------------------------------------
\newtcolorbox{descricao}{enhanced,breakable,boxrule=0.5pt,arc=1mm,
  left=2.5mm,right=2.5mm,top=1.6mm,bottom=2mm,
  colback=panel,colframe=plumAccent,colbacktitle=plumDeep,coltitle=white,
  fonttitle=\bfseries\small,title={Descrição}}

\newtcolorbox{funcao}{enhanced,breakable,boxrule=0.5pt,arc=1mm,
  left=2.5mm,right=2.5mm,top=1.6mm,bottom=2mm,
  colback=panelAlt,colframe=plumAccent,colbacktitle=plumDeep,coltitle=white,
  fonttitle=\bfseries\small,title={Função}}

\newtcolorbox{definicao}{enhanced,breakable,boxrule=0.5pt,arc=1mm,
  left=2.5mm,right=2.5mm,top=1.6mm,bottom=2mm,
  colback=panelAlt,colframe=plumAccent,colbacktitle=plumDeep,coltitle=white,
  fonttitle=\bfseries\small,title={Definição}}

\newtcolorbox{conceito}{enhanced,breakable,boxrule=0.5pt,arc=1mm,
  left=2.5mm,right=2.5mm,top=1.6mm,bottom=2mm,
  colback=panelAlt,colframe=plumAccent,colbacktitle=plumDeep,coltitle=white,
  fonttitle=\bfseries\small,title={Conceito}}

\newtcolorbox{metodo}{enhanced,breakable,boxrule=0.5pt,arc=1mm,
  left=2.5mm,right=2.5mm,top=1.6mm,bottom=2mm,
  colback=panelAlt,colframe=plumAccent,colbacktitle=plumDeep,coltitle=white,
  fonttitle=\bfseries\small,title={Método}}

\newtcolorbox{resultado}{enhanced,breakable,boxrule=0.5pt,arc=1mm,
  left=2.5mm,right=2.5mm,top=1.6mm,bottom=2mm,
  colback=panelAlt,colframe=plumAccent,colbacktitle=plumDeep,coltitle=white,
  fonttitle=\bfseries\small,title={Resultado e Cálculo}}

\newtcolorbox{limite}{enhanced,breakable,boxrule=0.5pt,arc=1mm,
  left=2.5mm,right=2.5mm,top=1.6mm,bottom=2mm,
  colback=panelAlt,colframe=plumAccent,colbacktitle=plumDeep,coltitle=white,
  fonttitle=\bfseries\small,title={Impacto e Limites}}

\newtcolorbox{arquitetura}{enhanced,breakable,boxrule=0.5pt,arc=1mm,
  left=2.5mm,right=2.5mm,top=1.6mm,bottom=2mm,
  colback=panelAlt,colframe=slate,colbacktitle=plumDeep,coltitle=white,
  fonttitle=\bfseries\small,title={Arquitetura e Elementos Técnicos}}

\newtcolorbox{niveis}{enhanced,breakable,boxrule=0.5pt,arc=1mm,
  left=2.5mm,right=2.5mm,top=1.6mm,bottom=2mm,
  colback=panelAlt,colframe=slate,colbacktitle=plumDeep,coltitle=white,
  fonttitle=\bfseries\small,title={Leitura do Elemento — Intuição $\rightarrow$ Implementação $\rightarrow$ Formalização}}

% Roteiro pedagógico usado em cada módulo: situa quem está começando,
% define o vocabulário, conduz à leitura técnica e termina com uma pergunta
% que exige evidência. A sequência evita pressupor conhecimento prévio.
\newcommand{\roteirodidatico}[5]{%
  \begin{tcolorbox}[enhanced,breakable,boxrule=0.55pt,arc=1mm,
    left=2.5mm,right=2.5mm,top=1.6mm,bottom=2mm,
    colback=panel,colframe=plumAccent,colbacktitle=plumDeep,coltitle=white,
    fonttitle=\bfseries\small,title={Roteiro de aprendizagem: do primeiro contato à análise avançada}]
    \textbf{Ponto de partida.} #1\par\smallskip
    \textbf{Vocabulário essencial.} #2\par\smallskip
    \textbf{Percurso técnico.} #3\par\smallskip
    \textbf{Leitura avançada.} #4\par\smallskip
    \textbf{Questão de verificação.} #5
  \end{tcolorbox}\par\smallskip
}
\newcommand{\casoguiado}[2]{%
  \begin{tcolorbox}[enhanced,breakable,boxrule=0.45pt,arc=1mm,
    left=2.5mm,right=2.5mm,top=1.4mm,bottom=1.8mm,
    colback=panelAlt,colframe=steel,colbacktitle=plumDeep,coltitle=white,
    fonttitle=\bfseries\small,title={Exemplo guiado — #1}]
    #2
  \end{tcolorbox}\par\smallskip
}

\newtcolorbox{aviso}{enhanced,breakable,boxrule=0.5pt,arc=1mm,
  left=2.5mm,right=2.5mm,top=1.6mm,bottom=2mm,
  colback=panelAlt,colframe=plumAccent,colbacktitle=plumDeep,coltitle=white,
  fonttitle=\bfseries\small,title={Aviso / Anti-overclaim}}

% Processo = fluxo de trabalho em etapas numeradas
\newtcolorbox{processo}{enhanced,breakable,boxrule=0.5pt,arc=1mm,
  left=2.5mm,right=2.5mm,top=1.6mm,bottom=2mm,
  colback=panelAlt,colframe=plumAccent,colbacktitle=plumDeep,coltitle=white,
  fonttitle=\bfseries\small,title={Fluxo de Trabalho (Processo)}}

% Referência auditada (biblioteca do apêndice)
\newtcolorbox{referencia}{enhanced,breakable,boxrule=0.5pt,arc=1mm,
  left=2.5mm,right=2.5mm,top=1.6mm,bottom=2mm,
  colback=panelAlt,colframe=plumAccent,colbacktitle=plumDeep,coltitle=white,
  fonttitle=\bfseries\small,title={Referência auditada}}

% Texto claro em todos os painéis editoriais do tema escuro.
\tcbset{coltext=textmain,fontupper=\color{textmain}}

% ----------------------------------------------------------------------
%  FLUXOGRAMAS (tikz)
% ----------------------------------------------------------------------
\tikzset{
  fn/.style={draw=lineDark,rounded corners=2pt,fill=panel,align=center,text=textmain,
             inner sep=3pt,minimum height=6mm,font=\scriptsize,text width=3.3cm},
  fns/.style={draw=steel,rounded corners=2pt,fill=steelV,align=center,text=textmain,
             inner sep=3pt,minimum height=6mm,font=\scriptsize\bfseries,
             text width=3.3cm},
  fne/.style={draw=gold,rounded corners=9pt,fill=navy,align=center,
             inner sep=3pt,minimum height=6mm,font=\scriptsize\bfseries,
             text=white,text width=3.3cm},
  fde/.style={draw=amber,fill=amber!45!navy,align=center,text=textmain,
             inner sep=3pt,minimum height=6mm,font=\scriptsize,
             diamond,aspect=2.2,text width=3.3cm},
  seta/.style={-Stealth,draw=textmuted,thick},
}
\newcommand{\nodep}[1]{{\urlstyle{tt}\scriptsize\bfseries\path{#1}}}

\newenvironment{flowch}[2][Fluxograma]{%
  \par\smallskip
  {\small\sffamily\bfseries\color{steel}#1}\hfill
  {\scriptsize\color{textmuted}#2}\par\smallskip
  \begin{center}\begin{tikzpicture}[font=\scriptsize]
}{%
  \end{tikzpicture}\end{center}\par\smallskip}

% ----------------------------------------------------------------------
%  LISTAS INTERPRETATIVAS
% ----------------------------------------------------------------------
\newenvironment{interpreta}{%
  \begin{description}[leftmargin=1.2em,style=nextline,itemsep=1pt]}{%
  \end{description}}

% ----------------------------------------------------------------------
%  RÓTULOS DE CÓDIGO / CAMINHOS — tipos literais seguros dentro de boxes:
%  o pacote "url" permite quebra de linha em / _ . - e mantém os caracteres
%  literais (sem erro de math mode com "_" etc.).
% ----------------------------------------------------------------------
\usepackage{url}
\newcommand{\pth}[1]{{\urlstyle{tt}\small\path{#1}}}
\newcommand{\cod}[1]{{\urlstyle{tt}\small\path{#1}}}
\emergencystretch=3em
\setlength{\hbadness}{4000}

% ----------------------------------------------------------------------
%  Conteúdo (Sumário) com hiperlinks
% ----------------------------------------------------------------------
\usepackage[nottoc]{tocbibind}

% ----------------------------------------------------------------------
%  SUMÁRIO: reserva largura suficiente para números de página romanos
%  longos (ex.: "xxii"/"xviii"), que estouravam ~0,6 pt cada entrada.
% ----------------------------------------------------------------------
\makeatletter
\renewcommand{\@pnumwidth}{2.2em}
\renewcommand{\@tocrmarg}{2.8em}
\renewcommand{\cftpartpagefont}{\color{textmain}}
\renewcommand{\cftchappagefont}{\color{textmain}}
\renewcommand{\cftsecpagefont}{\color{textmuted}}
\renewcommand{\cftsubsecpagefont}{\color{textmuted}}
\makeatother
